\documentclass[a4paper,12pt]{article}
\usepackage{cite}
\usepackage[a4paper, margin=1.5in]{geometry}
\usepackage{amsmath,amssymb,amsfonts,amsthm}
\usepackage{algorithmic}
\usepackage{graphicx}
\usepackage{textcomp}
\usepackage{xcolor}
\usepackage{txfonts}
\usepackage{listings}
\usepackage{enumitem}
\usepackage{graphicx}
\usepackage{mathtools}
\usepackage{gensymb}
\usepackage{comment}
\usepackage[siunitx, american]{circuitikz}
\usepackage[breaklinks=true]{hyperref}
\usepackage{tkz-euclide} 
\usepackage{listings}
\usepackage{gvv}        
\usepackage{float}
\usepackage{amsmath}
%\def\inputGnumericTable{}                                 
\usepackage[latin1]{inputenc}     
\usepackage{xparse}
\usepackage{color}                                            
\usepackage{array}                                            
\usepackage{longtable}                                       
\usepackage{calc}                                             
\usepackage{multirow}
\usepackage{multicol}
\usepackage{hhline}                                           
\usepackage{ifthen}     
\usepackage{placeins}
\usepackage{lscape}
\usepackage{tabularx}
\usepackage{array}
\usepackage{float}
\usepackage{listings}
\usepackage{xcolor}
\usepackage{booktabs}
\usepackage{caption}
\usepackage{tikz}
\usetikzlibrary{automata, arrows.meta}

\lstset{
  language=Python,
  basicstyle=\ttfamily\small,
  keywordstyle=\color{orange}\bfseries,
  stringstyle=\color{green!60!black},
  commentstyle=\color{gray},
  identifierstyle=\color{blue!70!black},
  numbers=left,
  numberstyle=\tiny\color{gray},
  frame=single,
  breaklines=true,
  showstringspaces=false,
  xleftmargin=0.2cm,
  xrightmargin=0cm,
  framexleftmargin=0.2cm,
  backgroundcolor=\color{gray!5}
}


\newtheorem{theorem}{Theorem}[section]
\newtheorem{problem}{Problem}
\newtheorem{proposition}{Proposition}[section]
\newtheorem{lemma}{Lemma}[section]
\newtheorem{corollary}[theorem]{Corollary}
\newtheorem{example}{Example}[section]
\newtheorem{definition}[problem]{Definition}
\newcommand{\BEQA}{\begin{eqnarray}}
\newcommand{\EEQA}{\end{eqnarray}}
\theoremstyle{remark}
\renewcommand{\thefigure}{\arabic{figure}}

\begin{document}
\title{Discrete - 1.1.63}
\author{ee25Btech11063 (Yalamati Vejith) }
\maketitle
\renewcommand{\thefigure}{\theenumi}
\renewcommand{\thetable}{\theenumi}
\textbf{Question}\\
\text{Quadrature formula:} $\int_{0}^{1} f(x) \, dx \approx \frac{1}{12} [f(0) + 2bf(0.25) + 2f(0.5) + 2df(0.75) + f(1)]$. \\
If used as Simpson's $\frac{1}{3}$ rule: \hfill (2007 XE) \\

\noindent a) b = d = 1 \hfill b) b = d = 2 \hfill c) b = 2d = 1 \hfill d) b = 2d = 2
\\ \textbf{Solution :}\\
\textbf{Simpson's 1/3 Rule}\\
Simpson's $1/3$ rule is a numerical integration method that approximates the area under a curve by fitting parabolic arcs (second-order polynomials) through sets of three points. 

For a single interval $[x_0, x_2]$ with a midpoint $x_1$ and step size $h$:
\begin{equation}
\int_{x_0}^{x_2} f(x) \, dx \approx \frac{h}{3} \left[ f(x_0) + 4f(x_1) + f(x_2) \right]
\end{equation}

In general \text{Simpson's 1/3 Rule} over $n$ intervals (where $n$ must be even):
\begin{equation}
\int_{a}^{b} f(x) \, dx \approx \frac{h}{3} \left[ f(x_0) + 4 \sum_{i=1,3,5}^{n-1} f(x_i) + 2 \sum_{j=2,4,6}^{n-2} f(x_j) + f(x_n) \right]
\end{equation}
Where $h = \frac{b-a}{n}$ is the uniform step size and $x_0$ = a, $x_1$ = b.\\ \\
The problem provides a specific quadrature formula:
\begin{equation}
\int_{0}^{1} f(x) \, dx \approx \frac{1}{12} \left[ f(0) + 2bf(0.25) + 2f(0.5) + 2df(0.75) + f(1) \right]
\label{eq:given}
\end{equation}

Based on the function arguments $\{0, 0.25, 0.5, 0.75, 1\}$:
\begin{itemize}
    \item $a = 0, b = 1$
    \item $n = 4$ intervals
    \item $h = \frac{1 - 0}{4} = 0.25$
\end{itemize}

Substituting $h = 0.25$ into the general Simpson's Rule for $n=4$:
\begin{equation}
\int_{0}^{1} f(x) \, dx \approx \frac{0.25}{3} \left[ f(0) + 4f(0.25) + 2f(0.5) + 4f(0.75) + f(1) \right]
\end{equation}

Simplifying the scalar $\frac{0.25}{3} = \frac{1}{12}$, the standard internal weights are:
\begin{equation}
\text{Internal Weights} = [4, 2, 4]
\end{equation}

Comparing these to the coefficients inside the brackets of the equation (3) we get
\begin{itemize}
    \item Coefficient for $f(0.25)$: $2b = 4 \implies \mathbf{b = 2}$
    \item Coefficient for $f(0.75)$: $2d = 4 \implies \mathbf{d = 2}$
\end{itemize}
\textbf{Numerical Example}\\
Lets test the above relation with $f(x) = x$
\\ \textbf{ Exact Integral :}\\
\begin{equation}
\int_{0}^{1} x \, dx = \left[ \frac{x^2}{2} \right]_0^1 = \frac{1}{2} = 0.5
\end{equation}
\textbf{ Applying the Formula} \\
\begin{equation}
\int_{0}^{1} f(x) \, dx \approx \frac{0.25}{3} \left[ f(0) + 4f(0.25) + 2f(0.5) + 4f(0.75) + f(1) \right]
\end{equation}
\begin{itemize}
    \item $f(0) = 0$
    \item $f(0.25) = 0.25 = \frac{1}{4}$
    \item $f(0.5) = 0.5 = \frac{1}{2}$
    \item $f(0.75) = 0.75 = \frac{3}{4}$
    \item $f(1) = 1$
\end{itemize}
Substitute these into above equation  and simplifying  it will give :
\begin{align}
\frac{1}{12} \left[ 0 + 4\left(\frac{1}{4}\right) + 2\left(\frac{1}{2}\right) + 4\left(\frac{3}{4}\right) + 1 \right]  \\
\implies \frac{1}{12} \left[ 1 + 1 + 3 + 1 \right] \\
\implies \frac{6}{12} = \frac{1}{2}
\end{align}
\newpage
\textbf{Python plot}
\begin{figure}[H]
    \centering
    \includegraphics[width=1.25\columnwidth]{figs/01.png}
    \label{fig:placeholder}
\end{figure}
\end{document}
